Starting from a one-step ChainOfThought system to solve ARC-AGI tasks, GEPA was able to evolve an elaborate 5-step AI system, that performs the following steps:
\begin{enumerate}
    \item Induce Natural Language Rule from training examples using LLM.
    \item Generate a python program based on the induced rule using LLM.
    \item Execute the python program on training examples, identifying failures and gathering feedback text from its execution and evaluation.
    \item If there are any observed failures on the training examples, use an LLM to improve the python program with text feedback.
    \item Execute the final python program on test inputs, and return outputs.
\end{enumerate}

This 5-step system evolved by GEPA improved Gemini-2.5-Pro's performance on ARC-AGI-1 from 44\% to 49.5\%. We present the full system evolved by GEPA:
\begin{instructbox}[\incode{ARC-AGI Gemini-2.5-Pro agent}]
Base Agent:
\begin{lstlisting}[breakindent=0pt, breaklines=true, breakatwhitespace=true, frame=single, backgroundcolor=\color{gray!10}]
import dspy
from typing import List
import pydantic

MATRIX = List[List[int]]

class TrainingExample(pydantic.BaseModel):
    input: MATRIX
    output: MATRIX

class SolveTaskSignature(dspy.Signature):
    training_examples: List[TrainingExample] = dspy.InputField(description="Input and output examples demonstrating the task to be performed.")
    test_inputs: List[MATRIX] = dspy.InputField(description="Input matrices to be solved following the task described in the training examples.")
    test_outputs: List[MATRIX] = dspy.OutputField(description="Output matrices corresponding to the test inputs.")

program = dspy.ChainOfThought(SolveTaskSignature)
\end{lstlisting}

GEPA Evolved Agent:
\begin{lstlisting}[breakindent=0pt, breaklines=true, breakatwhitespace=true, frame=single, backgroundcolor=\color{blue!10}]
import dspy
from typing import List, Tuple, Optional
import pydantic
import copy
import traceback

# Define type aliases and pydantic models for clarity and structure.
MATRIX = List[List[int]]

class TrainingExample(pydantic.BaseModel):
    input: MATRIX
    output: MATRIX

# --- Signatures ---

class HypothesizeRule(dspy.Signature):
    """
    Analyze the provided input/output matrix pairs from the Abstraction and Reasoning Corpus (ARC).
    Deduce the single, underlying transformation rule that converts each input matrix to its corresponding output matrix.
    Describe this rule in clear, step-by-step, unambiguous English. Focus on the logic, not Python code.

    **Successful Strategies to Consider:**
    - **Start Simple:** First, check for simple rules. Is there a global transformation (e.g., rotation, reflection)? Is the output a subgrid of the input? Is a single color being replaced?
    - **Look for Separators:** Check if the grid is partitioned by separator lines (e.g., rows or columns that are all one color, usually black). The transformation might be applied to each section independently.
    - **Identify Objects:** Group contiguous non-background pixels into objects. Analyze how these objects are created, destroyed, or modified. Consider their properties: color, shape, size, position.
    - **Find Marker Points:** Look for special 'marker' pixels (e.g., a uniquely colored pixel) that might define the geometry of an operation, like the corners of a shape to be drawn.
    - **Relate Input to Output:** How do the properties of the input grid (dimensions, colors, object counts) relate to the output grid?
    - **Avoid Over-complication:** Propose the simplest rule that explains ALL training examples. Do not suggest overly complex mathematical or recursive patterns unless absolutely necessary and supported by every example.
    """
    training_examples: List[TrainingExample] = dspy.InputField(desc="A list of input/output pairs demonstrating the transformation rule.")
    rule_description: str = dspy.OutputField(desc="A step-by-step English description of the transformation rule.")

class ImplementRule(dspy.Signature):
    """
    You are an expert programmer. Your task is to write a single, self-contained Python function based on a provided rule description and example pairs.

    **Function Requirements:**
    - The function must be named `transform_matrix`.
    - It must accept one argument: `matrix`, which is a list of lists of integers (representing the input grid).
    - It must return a list of lists of integers (representing the transformed output grid).
    - The function should not use any external libraries except for `copy` if needed (e.g., `import copy; new_matrix = copy.deepcopy(matrix)`).
    - Your output must be ONLY the Python code for the function. Do not include any explanations, comments outside the function, or markdown formatting like ```python.

    **Successful Strategies to Consider:**
    - **Robust Parsing:** When rules involve sub-regions (e.g., quadrants, objects), write code that can robustly find their boundaries. Don't assume boundaries are perfectly aligned or can be found by checking just the first row/column.
    - **Iterative Processes:** Some rules are applied repeatedly until the grid no longer changes. Consider using a `while` loop that continues as long as modifications are being made in a pass.
    - **Consistent Output Shape:** When constructing the output grid row by row (e.g., from sections), ensure each generated row has the same length to avoid creating a staggered/jagged matrix.
    - **Handling Edge Cases:** Ensure your code handles empty matrices or matrices with unexpected dimensions gracefully.

    **Example of a Correctly Formatted Output:**
    def transform_matrix(matrix: list[list[int]]) -> list[list[int]]:
        # Your implementation here
        # For example, find the most frequent color and fill the grid
        from collections import Counter
        import itertools
        
        if not matrix or not matrix[0]:
            return []
            
        counts = Counter(itertools.chain.from_iterable(matrix))
        if counts:
            # Handle ties by picking the smaller number value
            most_common_color = sorted(counts.items(), key=lambda item: (-item[1], item[0]))[0][0]
        else:
            return []

        height = len(matrix)
        width = len(matrix[0])
        
        return [[most_common_color for _ in range(width)] for _ in range(height)]
    """
    rule_description: str = dspy.InputField(desc="The English description of the rule to implement.")
    training_examples: List[TrainingExample] = dspy.InputField(desc="A list of input/output pairs to use as a reference for implementation.")
    python_function: str = dspy.OutputField(desc="A string containing a single Python function `transform_matrix` that implements the rule.")

class RefineCode(dspy.Signature):
    """
    You are a senior programmer debugging a Python function. You will be given a rule description, training examples, the buggy Python code, and feedback on why it failed.
    Your task is to fix the code so that it correctly implements the rule and passes the training examples.
    The refined function must adhere to all the original requirements.

    **Function Requirements:**
    - The function must be named `transform_matrix`.
    - It must accept one argument: `matrix`, which is a list of lists of integers.
    - It must return a list of lists of integers.
    - The function should not use any external libraries except for `copy`.
    - Your output must be ONLY the Python code for the function. Do not include any explanations or markdown formatting.
    """
    rule_description: str = dspy.InputField(desc="The English description of the rule to implement.")
    training_examples: List[TrainingExample] = dspy.InputField(desc="A list of input/output pairs the function must correctly handle.")
    buggy_code: str = dspy.InputField(desc="The Python code that failed verification.")
    feedback: str = dspy.InputField(desc="A description of the error or the discrepancy between the function's output and the expected output for a training example.")
    refined_python_function: str = dspy.OutputField(desc="A string containing the corrected, single Python function `transform_matrix`.")

# --- Custom Module with Self-Correction ---

class ARCSolver(dspy.Module):
    """A module that solves ARC tasks by hypothesizing a rule, generating code, and refining it through verification."""
    def __init__(self, max_retries=2):
        super().__init__()
        self.max_retries = max_retries
        self.rule_hypothesizer = dspy.ChainOfThought(HypothesizeRule)
        self.code_implementer = dspy.Predict(ImplementRule)
        self.code_refiner = dspy.Predict(RefineCode)

    def _execute_and_verify(self, python_code: str, examples: List[TrainingExample]) -> Tuple[bool, Optional[str]]:
        """
        Executes the generated Python code and verifies its correctness against training examples.
        Returns a tuple: (is_correct, feedback_message).
        """
        local_scope = {}
        try:
            exec(python_code, globals(), local_scope)
            transform_func = local_scope.get('transform_matrix')
            if not callable(transform_func):
                return False, "The generated code does not define a callable function named 'transform_matrix'."
        except Exception as e:
            return False, f"Code compilation failed with an error: {e}\n{traceback.format_exc()}"

        for i, example_obj in enumerate(examples):
            # Robustly handle both Pydantic models and dicts, as DSPy can convert them.
            if isinstance(example_obj, dict):
                try:
                    example = TrainingExample(**example_obj)
                except pydantic.ValidationError as e:
                    return False, f"Failed to parse training example {i} from dict: {e}"
            else:
                example = example_obj

            try:
                input_copy = copy.deepcopy(example.input)
                predicted_output = transform_func(input_copy)
                if predicted_output != example.output:
                    feedback = (
                        f"Verification failed on training example {i}.\n"
                        f"Input:\n{example.input}\n"
                        f"Expected Output:\n{example.output}\n"
                        f"Actual Output:\n{predicted_output}"
                    )
                    return False, feedback
            except Exception as e:
                feedback = (
                    f"An exception occurred during execution on training example {i}: {e}\n"
                    f"Input was:\n{example.input}\n"
                    f"{traceback.format_exc()}"
                )
                return False, feedback
        
        return True, None

    def forward(self, training_examples: List[TrainingExample], test_inputs: List[MATRIX]):
        # Step 1: Generate an English description of the transformation rule.
        hypothesis = self.rule_hypothesizer(training_examples=training_examples)
        
        # Step 2: Generate the initial Python code.
        prediction = self.code_implementer(
            rule_description=hypothesis.rule_description,
            training_examples=training_examples
        )
        python_code = prediction.python_function

        # Step 3: Verify and refine the code in a loop.
        for _ in range(self.max_retries):
            is_correct, feedback = self._execute_and_verify(python_code, training_examples)
            if is_correct:
                break
            
            refinement = self.code_refiner(
                rule_description=hypothesis.rule_description,
                training_examples=training_examples,
                buggy_code=python_code,
                feedback=feedback
            )
            python_code = refinement.refined_python_function
        
        # Step 4: Execute the final code on test inputs with robust fallbacks.
        fallback_outputs = [copy.deepcopy(matrix) for matrix in test_inputs]
        local_scope = {}
        try:
            exec(python_code, globals(), local_scope)
            transform_func = local_scope.get('transform_matrix')
            if not callable(transform_func):
                return dspy.Prediction(test_outputs=fallback_outputs)

            solved_outputs = []
            for test_matrix in test_inputs:
                try:
                    input_copy = copy.deepcopy(test_matrix)
                    result = transform_func(input_copy)
                    solved_outputs.append(result)
                except Exception:
                    solved_outputs.append(copy.deepcopy(test_matrix))
            
            return dspy.Prediction(test_outputs=solved_outputs)

        except Exception:
            return dspy.Prediction(test_outputs=fallback_outputs)

# The overall task signature, defining the final inputs and outputs of the program.
class SolveTaskSignature(dspy.Signature):
    """
    Given a set of training examples demonstrating a matrix transformation,
    apply the same transformation to a new set of test matrices.
    """
    training_examples: List[TrainingExample] = dspy.InputField(description="Input and output examples demonstrating the task to be performed.")
    test_inputs: List[MATRIX] = dspy.InputField(description="Input matrices to be solved following the task described in the training examples.")
    test_outputs: List[MATRIX] = dspy.OutputField(description="Output matrices corresponding to the test inputs.")

# The final program is an instance of our new, more robust module.
program = ARCSolver()
\end{lstlisting}
\end{instructbox}